\section{增长叙事必须落到EPS}

\begin{exhibitbox}[Exhibit 10：已启动核心的增长盈利桥]
\begin{tabularx}{\textwidth}{L{2.2cm}L{3.0cm}L{2.6cm}L{2.4cm}X}
\toprule
\textbf{公司} & \textbf{增长代理} & \textbf{2026E收入/净利} & \textbf{现价隐含} & \textbf{估值信用} \\
\midrule
中兴通讯 & 算力收入占比27\%、项目交付 & 1500/71亿元 & 27.4x PE & 期权计入合并PE，不单独拔高算力倍数 \\
江波龙 & H1预告、存储价格、企业级SSD & 440/200亿元 & 12.4x PE & 盈利信用，但给予周期和现金流折价 \\
恒瑞医药 & 创新药销售、BD、获批/临床 & 360/94亿元 & 39.3x PE & SOTP盈利信用，不对仿制药套创新药倍数 \\
工业富联 & AI服务器收入、800G交换机、H1预告 & 14804/560亿元 & 23.5x PE & 预告后合并PE，不单独拔高AI分部倍数 \\
\bottomrule
\end{tabularx}
\end{exhibitbox}

\section{中兴通讯：算力份额不是利润的同义词}

中兴的增长业务采用算力收入占比作为公开代理：按2026E收入1500亿元、算力占比27\%估算，增长业务收入约405亿元，连接与其他基本盘约1095亿元。由于公司未披露算力业务毛利率、设备数量、ASP和客户分配，本报告不对405亿元收入单独赋予高增长PS或PE，只在合并PE中给予期权信用。

当前40.53元若按33倍基准PE，隐含EPS约1.23元、净利润约59亿元；公开一致预期71亿元已经高于这一门槛。真正决定48.75元目标能否兑现的，是H2利润修复、算力业务毛利率和经营现金流，而不是算力收入占比继续上升本身。

\section{江波龙：官方H1预告关闭了上半年利润桥}

2025年收入227.7亿元、归母净利润14.2亿元构成基本盘。2026E收入440亿元意味着增量约212亿元；H1官方预告已经贡献收入220--250亿元和净利润92--110亿元。基准全年净利润200亿元要求H2约99亿元，略低于H1中值，属于“高景气延续但不继续加速”的假设。

当前价若按19倍PE，隐含全年净利润约131亿元，仅要求H2在H1中值基础上再贡献约30亿元。这说明现价确实保留了周期折价。但如果现金流持续为负、低价库存红利结束或NAND价格反转，合理倍数会快速收缩，520元熊市价值并非极端情景。

\section{恒瑞医药：创新药与成熟业务必须分部估值}

恒瑞2026E收入约360亿元。按华泰对创新药占比67\%的判断，创新药与BD约241亿元，成熟产品和其他收入约119亿元。创新药部分可用风险调整DCF和管线价值，成熟业务只能用较低倍数；这也是采用SOTP而不是给全公司统一高增长PE的原因。

现价对应2026E PE约39倍，并不属于绝对低估。剩余空间来自创新药销售30\%以上增长、BD现金流和新产品获批，而不是简单的估值均值回归。若创新药增长低于20\%、EPS低于1.25元或核心临床里程碑延后，SOTP必须下修。

\section{工业富联：H1预告关闭利润桥，但H2仍需321亿元}

工业富联2025年收入9028.9亿元是基本盘，华泰预告后预测2026E收入14,803.7亿元、净利560.1亿元、EPS 2.82元。H1净利中值239亿元意味着H2需要约321.1亿元才能完成华泰基准。公司未披露AI服务器具体台数、单机ASP、客户分配和分部毛利率，因此本报告只使用AI服务器收入增速、800G以上交换机出货和下一代平台量产作为代理，不对AI分部另给高增长倍数。

现价66.27元对应华泰EPS约23.5倍PE。AStock基准30倍得到84.60元，最终多锚目标82.29元。若H2净利低于321亿元、AI服务器增速跌破100\%或毛利率低于7\%，必须下修基准倍数。

\section{四种执行纪律}

\begin{itemize}
  \item \textbf{中兴通讯}：启动后回踩确认。39--40元附近观察承接，站稳43--44元且H2利润修复再升级。
  \item \textbf{江波龙}：只做高波动回撤。560--588元为第一观察区，跌破520元且现金流恶化则取消增长信用。
  \item \textbf{恒瑞医药}：趋势回撤配置。50--55元分批，57.5--61元突破后必须由创新药销售和里程碑继续验证。
  \item \textbf{工业富联}：业绩兑现回撤。62--66元观察承接，突破72--76元必须由H2量产、净利和毛利率继续验证。
\end{itemize}
